\chapter*{Abstract}
\addcontentsline{toc}{chapter}{Abstract}
This thesis develops a causal theory for Lamb waves generated by thermoelastic
laser sources in isotropic plates. Conventional frequency-domain calculations
give the response to a harmonic source that has acted indefinitely; laboratory
sources are switched on, pulsed for finitely many cycles, scanned, or patterned.
The central problem is therefore to determine when a declared measurement has
approached the forced periodic response and how this differs from first arrival.

Starting from linear elasticity, the plate eigenproblem, Rayleigh--Lamb equations,
eigenfields, energy flux, and normalization are derived. Optical absorption and
heat diffusion enter through thermal strain, and a weak-work projection gives
the physical modal force. Causal Green functions separate forced periodic and
homogeneous transient parts. Propagating stationary points, cutoffs, and ZGV
points are treated as distinct regimes. Existing scalar-overlap transients are
retained only as preliminary software checks, not physical modal purity or
absolute laser-coupling predictions.
